\documentclass[11pt,a4paper]{article}
\usepackage[margin=1in]{geometry}
\usepackage{enumitem}
\usepackage{hyperref}
\hypersetup{colorlinks=false}
\pagestyle{empty}

\begin{document}

\noindent
\today

\bigskip

\noindent
The Editor-in-Chief\\
\textit{International Journal of Forecasting}\\
International Institute of Forecasters (Elsevier)

\bigskip

\noindent
Dear Editor,

\medskip

I am pleased to submit ``When Does Better Volatility Measurement Improve Allocation? A Cross-Asset Complexity Ceiling'' for consideration in the \textit{International Journal of Forecasting}. The paper asks a deliberately narrow forecasting question---when does an improvement in volatility \emph{measurement} actually improve a downstream portfolio \emph{allocation}?---and answers it with a pre-specified out-of-sample evaluation carried through from forecast accuracy to a volatility-targeting decision.

The central finding is a \emph{complexity ceiling}. Moving from a symmetric GARCH(1,1) to an asymmetric GJR-GARCH(1,1), and selecting on the sign and significance of the asymmetry parameter, earns statistically detectable out-of-sample forecast gains only in a restricted domain: a homogeneous equity universe with a significant asymmetry. Even where measurement improves at the forecast level, the improvement largely fails to reach the allocation level, where the dominant and universal benefit of volatility targeting---drawdown reduction---is governed by an asset's base volatility level rather than by its asymmetric structure. The economic content of the asymmetry (balance-sheet leverage for equities, regime-dependent flight-to-quality for gold, coupon-dominated pricing for bonds) is an interpretable classification device but does not travel across asset classes.

The manuscript is well suited to this journal for three reasons. First, its contribution is squarely a forecast-evaluation result: it uses the QLIKE loss, Diebold--Mariano inference, and multiple-testing discipline to establish when added model complexity does---and does not---earn its keep out of sample. Second, it is an honest, decision-relevant negative result of the kind the \textit{International Journal of Forecasting} regularly publishes, showing that in-sample refinement does not reliably survive to the decision level. Third, credibility rests on a pre-specified evaluation protocol and a self-contained replication package (data and code) rather than on an ex-post specification search; a reproducibility checker verifies every reported table number against its canonical source.

The manuscript has not been submitted elsewhere and is not under consideration at any other journal. I suggest the following potential referees, whose work is central to the forecast-evaluation methods the paper employs:

\begin{enumerate}[nosep]
\item \textbf{Andrew J.\ Patton} (Duke University) --- Volatility forecast comparison under imperfect proxies and the QLIKE loss.
\item \textbf{Peter R.\ Hansen} (University of North Carolina at Chapel Hill) --- Model confidence sets and GARCH forecast comparison.
\item \textbf{Kevin Sheppard} (University of Oxford) --- Volatility modelling and out-of-sample evaluation.
\end{enumerate}

\noindent Thank you for considering my work. I look forward to hearing from you.

\bigskip

\noindent
Sincerely,

\bigskip

\noindent
Yi-Hao Lai\\
Department of Finance, Da-Yeh University\\
Email: yhlai@mail.dyu.edu.tw

\end{document}
